\documentclass[10pt]{article}
\usepackage[margin=0.7in]{geometry}
\usepackage[T1]{fontenc}
\usepackage[utf8]{inputenc}
\usepackage{lmodern}
\usepackage{amsmath,amssymb,microtype}
\usepackage[hidelinks]{hyperref}
\setlength{\parindent}{0pt}
\setlength{\parskip}{0.3em}
\setlength{\abovedisplayskip}{5pt}
\setlength{\belowdisplayskip}{5pt}
\pagestyle{empty}
\begin{document}
\begin{center}\large\bfseries Updated Climate-Zone Heat Model\end{center}

\textbf{1. Full current model.}
Let $i$ index ZIP codes, $d$ days, $k(i)$ SES clusters, and $\mathrm{cz}[i]$
climate zones. Retain the existing standardized heat $H_{id}$, ZIP mean
$\bar H_i$, acute heat $x_{id}=H_{id}-\bar H_i$, and population exposure $n_{id}$.
The updated model for ED counts is
\begin{align}
Y_{id}\mid\theta&\sim\operatorname{NB}(n_{id}\mu_{id},r),\qquad r>0,\nonumber\\
\log\mu_{id}&=\alpha+\gamma_{k(i)}+\beta_{\mathrm{chronic}}\bar H_i
+\beta_{\mathrm{acute},k(i)}(H_{id}-\bar H_i)
+\beta_{m(d)}+\beta_{w(d)}\nonumber\\
&\quad+f_{\mathrm{doy}}(d)
+f_{\mathrm{acute},\mathrm{cz}[i]}(H_{id}-\bar H_i)+u_i,\label{eq:model}\\
f_{\mathrm{acute},\mathrm{cz}[i]}(\cdot)
&=f_{\mathrm{shared}}(\cdot)+f_{\mathrm{dev},\mathrm{cz}[i]}(\cdot),\nonumber\\
\beta_{\mathrm{acute},k}&=\beta_{\mathrm{acute}}+\delta_k,\qquad
u_i\mid\tau_{\mathrm{ZIP}}\overset{\mathrm{iid}}{\sim}N(0,\tau_{\mathrm{ZIP}}^{-1}).\nonumber
\end{align}
Here, $u_i$ is the ZIP-specific random intercept, allowing each ZIP to have
its own underlying ED rate after accounting for the other predictors.
The precision $\tau_{\mathrm{ZIP}}$ is its inverse variance; iid denotes
independent, identically distributed intercepts conditional on that precision.
The terms $m(d)$ and $w(d)$ denote month and weekday; $f_{\mathrm{doy}}$ is the
seasonal RW2 effect. Reference-cluster constraints remain
$\gamma_{k_0}=\delta_{k_0}=0$, with $k_0=6$.
The SES-specific heat slopes are linear; the RW2 heat curves supply the
nonlinear response.

Acute-heat contrasts include the SES-specific linear slope and both heat
curves. Comparing $x=0$ (the ZIP's mean heat) with $x=1$ (one county-wide
heat SD above that mean),
\[
L_i=\beta_{\mathrm{acute},k(i)}
+f_{\mathrm{acute},\mathrm{cz}[i]}(1)
-f_{\mathrm{acute},\mathrm{cz}[i]}(0),\qquad
\mathrm{IRR}_{1\mathrm{SD},i}=e^{L_i},\qquad
p_i=\Pr(L_i>0\mid Y).
\]
INLA linear-combination marginals provide medians, 95\% credible intervals
and probabilities for the 42 observed SES-cluster--zone combinations,
including missing zones. Curves are evaluated at the nearest nodes on the
retained 50-node grid. Excess-rate uncertainty was estimated using 2,000
joint posterior draws that respected the fitted model constraints.

\textbf{2. ZIP baseline term.}
The model includes \texttt{f(zip\_id, model = "iid")}.
The ZIP-level SD, $\sigma_{\mathrm{ZIP}}=\tau_{\mathrm{ZIP}}^{-1/2}$,
is summarized from the transformed precision marginal, giving a posterior
mean of 0.278 and 95\% credible interval 0.231--0.321 on the log-rate scale.
Fit and zero-fraction diagnostics are shown under Overdispersion.
Conditional baseline cluster effects are reported as $\exp(\gamma_k)$ with 95\% credible
intervals and reference cluster 6 assigned IRR 1. A matched comparison
with and without the ZIP term is not implemented, so changes cannot be
attributed to adding that term alone; the likelihood comparison below
retains it in both models.

\textbf{3. Deviation structure.}
Independent replicated climate-zone RW2 deviations are replaced by grouped
deviations using \texttt{group = climate\_zone\_group} and
\texttt{control.group = list(model = "exchangeable")}.
Each of the five zones (20, 22, 28, 29 and 30) has an adjustment curve
$f_{\mathrm{dev},\mathrm{cz}}$ to the shared heat curve. These curves retain
a common curvature precision $\tau_{\mathrm{dev}}$, with between-zone
correlation matrix
\[
R(\rho)=
\begin{pmatrix}
1&\rho&\rho&\rho&\rho\\
\rho&1&\rho&\rho&\rho\\
\rho&\rho&1&\rho&\rho\\
\rho&\rho&\rho&1&\rho\\
\rho&\rho&\rho&\rho&1
\end{pmatrix},
\qquad -\tfrac14<\rho<1.
\]
Each row and column represents one zone. Every pair shares correlation
$\rho$ between its adjustment curves, after accounting for the shared
curve. The bounds ensure a valid, nonsingular five-zone correlation matrix.
For ZIPs without a zone assignment, the zone deviation is zero; the other
predictor terms remain.

\textbf{4. Priors.}
The seasonal RW2, shared-heat RW2, zone-deviation RW2 and ZIP-intercept
precisions each have a PC prior. For component $j$, write its scale as
$\sigma_j=\tau_j^{-1/2}=1/\sqrt{\tau_j}$; the baseline calibration is
\[
\Pr(\sigma_j>1)=0.01,\qquad
j\in\{\mathrm{doy},\mathrm{shared},\mathrm{dev},\mathrm{ZIP}\}.
\]
Here $j$ labels the four components; each has its own estimated precision.
For ZIP intercepts, $\sigma_j$ is their SD. For scaled RW2 terms, it controls
the penalized curve variation: these priors discourage bending, while
straight-line trends remain unpenalized by the RW2 prior.
The correlation prior favors $\rho=0$, with prior probability $0.5$ on
negative correlations and $0.5$ on positive correlations. Conditional on
either sign, the distance $d_\rho=\sqrt{-\log\det R(\rho)}$ from zero
correlation has an exponential distribution, calibrated so that
$\Pr(\rho>0.5)=0.1$. This prior is converted to INLA's computation scale
without changing its probabilities.
NB size $r$ uses \texttt{pc.mgamma} with INLA's default rate parameter 7,
shrinking toward the Poisson limit. The 39 non-intercept fixed coefficients
retain independent $N(0,1000)$ priors, where 1000 is the variance;
the overall intercept $\alpha$ remains flat.

For precision-prior sensitivity, the four precision priors were varied
jointly: stronger PC priors use
$\Pr(\tau_j^{-1/2}>0.5)=0.01$, weaker PC priors use
$\Pr(\tau_j^{-1/2}>2)=0.01$, and a comparison uses
$\tau_j\sim\operatorname{Gamma}(1,5\times10^{-5})$ with shape--rate
parameters. Correlation, NB-size and fixed-coefficient priors remain at
baseline. These fit-only comparisons use heat-effect marginals,
conditional SES-cluster IRRs, DIC and WAIC, without using additional
excess-rate draws.

For coefficient-variance sensitivity, separate NB fits changed only the
Normal prior variance for the 39
non-intercept fixed coefficients from 1000 to 100 or 10, retaining zero prior
means, the flat intercept and all PC priors. Variance 100 was optimized
again from its fitted solution; this retained fit converged without
likelihood-recovery warnings. The diagnostics are summarized below.
\begin{center}
\small
\setlength{\tabcolsep}{5pt}
\renewcommand{\arraystretch}{1.15}
\begin{tabular}{lrrrrr}
\hline
\shortstack{Normal prior\\variance} & DIC & WAIC &
\shortstack{Predicted\\zeros (\%)} &
\shortstack{Max. IRR median\\change (\%)} &
\shortstack{Max. probability\\change (pp)}\\
\hline
1000 (baseline) & 65,207.33 & 65,254.86 & 86.98 & 0.00 & 0.00\\
100 (rerun) & 65,207.28 & 65,254.33 & 86.98 & 0.77 & 1.14\\
10 & 65,206.02 & 65,253.74 & 86.99 & 0.52 & 0.81\\
\hline
\end{tabular}
\end{center}
Changes are relative to baseline across the 42 complete heat contrasts;
pp denotes percentage points. The retained variance-100 fit agreed with baseline on
whether each of the 42 heat intervals excluded 1. At variance 10, one borderline
interval (SES cluster 14, zone 29) changed from including to excluding 1.
These fit-only checks used INLA marginals without additional excess-rate draws; the
variance-1000 NB baseline was retained.

\textbf{5. Overdispersion.}
The negative-binomial likelihood allows daily ED counts to vary more than
a Poisson model would allow around the same predicted mean.
Its conditional mean is $M_{id}=n_{id}\mu_{id}$ and variance is
$M_{id}+M_{id}^2/r$; as $r\to\infty$, it approaches Poisson.
A matched Poisson fit retained the ZIP intercept, predictor terms,
population exposures and common priors on the same 54,910 ZIP-days.
Observed zeros numbered 47,898 (87.23\%). The comparison supports retention
of the NB baseline; smaller DIC and WAIC are preferred.
\begin{center}
\small
\renewcommand{\arraystretch}{1.15}
\begin{tabular}{lrrr}
\hline
Count likelihood & DIC & WAIC & Predicted zeros (\%)\\
\hline
NB baseline & 65,207.33 & 65,254.86 & 86.98\\
Matched Poisson & 91,394.49 & 94,721.52 & 74.16\\
\hline
\end{tabular}
\end{center}
Zero probabilities use posterior mean
fitted rates and, for NB, posterior mean size: these are in-sample
plug-in diagnostics. Checks also cover clusters, ZIPs, zones and months.

\textbf{6. Suppressed counts.}
In \texttt{suppressed\_data.R}, daily suppressed cells are initially
imputed by uniform integer draws from 1--11 (seed 1), then adjusted
toward monthly totals using proportional allocation and largest-remainder
rounding. Original daily zeros are preserved before the monthly-allocation
branches. The special one-visit override for all-zero ZIP-months is not
used: the seven ZIP-months with a suppressed monthly total and all-zero
daily records remain zero and are exceptions to reconciliation. The
initial 1--11 range is not reapplied as a per-day upper bound after
reconciliation.

\textbf{7. Data source.}
The ED inputs are the dummy files
\path{Dummy_20241005_TotalERDaily.csv} and
\path{Dummy_20241005_TotalERMonthly.csv}.
Zero heat-index values denote unavailable heat and are excluded from
fitting; genuine zero ED counts are retained. Catalina (90704) is kept:
the coverage report documents usable heat on 176 of 184 days (95.7\%).
Its missing climate-zone assignment is handled by retaining the shared
heat curve, SES-specific slope and ZIP intercept without a zone deviation.
The mainland map extent restricts display, not model eligibility.

\textbf{8. Excess-rate reference.}
Unchanged: excess rates compare observed heat with the SES cluster's monthly mean heat.

\textbf{9. Figures.}
Cluster outlines: yes. Cluster numbering: yes. Forest plots: yes.
Acute-heat IRRs per SD with $^{\circ}$F labels: yes.
Baseline cluster IRRs labeled separately: yes.

\textbf{Other.}
Excess-rate calculations include the complete heat response. The ZIP
intercept and other non-heat terms are held fixed within each contrast.
Signed rate differences are averaged with equal weight for each modeled
ZIP-day and reported per 10,000 population per day; positive and negative
contributions are both retained. Conditional sampling variances were
checked against those saved by INLA. Zero-fraction diagnostics use
in-sample plug-in predictions rather than fully integrated
posterior-predictive tests.

\end{document}
